\documentclass[12pt,letterpaper]{article}

\usepackage[utf8]{inputenc}
\usepackage[spanish]{babel}
\usepackage{graphicx}
\usepackage{geometry}
\usepackage{amsmath}
\usepackage{amssymb}
\usepackage{float}
\usepackage{caption}
\usepackage{subcaption}
\usepackage{setspace}
\usepackage{csquotes}
\usepackage{booktabs}
\usepackage{titling}
\usepackage{siunitx}
\usepackage{url}
\usepackage{tikz}
\usetikzlibrary{positioning, arrows.meta}
\usepackage[backend=biber,style=ieee]{biblatex}
\addbibresource{references.bib}

\geometry{left=2.5cm, right=2.5cm, top=2.5cm, bottom=2.5cm}

\graphicspath{{figuras/}}

% --- Figura con placeholder automático: compila con o sin la imagen ----------
\usepackage{xcolor}
\newcommand{\figph}[3]{%
  \begin{figure}[htbp]
    \centering
    \IfFileExists{figuras/#1}{%
      \includegraphics[width=0.70\textwidth]{#1}%
    }{%
      \fcolorbox{gray}{gray!8}{\begin{minipage}[c][0.15\textheight][c]{0.68\textwidth}%
        \centering\footnotesize\color{gray!55!black}%
        [\,PLACEHOLDER\,]\\[4pt] \texttt{\detokenize{figuras/#1}}%
      \end{minipage}}%
    }
    \caption{#2}
    \label{fig:#3}
  \end{figure}%
}

% =============================================================================
% Reporte de Avance — Cuarto Semestre. Estado: SEMIFINAL.
% =============================================================================

\title{Análisis de Tejidos mediante Tomografía Fotoacústica e IA\\
El Sistema de Adquisición de Datos (DAQ) de Alta Velocidad:\\
Integración y Validación con ADC Físico AD9226}
\author{Eduardo Santos Mena \\ Asesores: Dr. Samuel Pérez Huerta, Dr. Augusto David Ariza Flores \\
\small Unidad Académica de Ciencia y Tecnología de la Luz y la Materia, UAZ}
\date{Reporte de Avance — Cuarto Semestre: enero–junio 2026}

\begin{document}

\maketitle

\begin{abstract}
\noindent
El avance central del periodo es la consolidación y validación de un \textbf{sistema
de adquisición de datos (DAQ) de alta velocidad y bajo costo} para fotoacústica.
Dando continuidad al periodo anterior —núcleo digital validado con señales sintéticas—,
se completó la integración física del conversor analógico-digital (AD9226) con la FPGA
y se certificó, con señal real, que la cadena captura, almacena y transmite
\textbf{sin pérdida de datos} a \textbf{54~MS/s con 12 bits} (LSB 2.44~mV). Llevar el
ciclo de lectura completo a esa resolución y tasa —de 8 a 12 bits y de 27 a 54~MS/s—
multiplica el volumen de datos por ventana, con consecuencias concretas: mayor
ocupación de BRAM, una transmisión UART más larga (menor tasa de ráfagas) y una
sincronización más exigente de la ventana de captura. En la detección se observó, en
el osciloscopio, una señal en el transductor con apariencia de \textit{ringing}
acústico ($\sim$2.5~MHz); sin embargo, tres controles la identifican como acoplamiento
electromagnético y, al conectar el transductor al ADC, la señal se perdía —solo
persistía en el osciloscopio—, lo que evidencia el problema de acoplamiento de
impedancias. Aunque la validación fotoacústica completa (transductor excitado y
sensando) no se concretó, diseñar y programar la FPGA como DAQ a estas frecuencias y
con materiales de bajo costo constituye un \textbf{hito}: deja una plataforma de
adquisición robusta y acota la detección a un problema bien definido de front-end
analógico.
\end{abstract}

\section{Introducción}

La tomografía fotoacústica (PAT) combina la especificidad molecular de la excitación
óptica con la resolución de la detección ultrasónica, sin radiación
ionizante~\cite{lin_emerging_2022}. Este trabajo se enmarca en un proyecto de
doctorado (LUMAT--UAZ) cuyo objetivo general es desarrollar y evaluar un sistema
integrado de PAT y \textit{deep learning} para la detección temprana de la artritis
reumatoide (AR)~\cite{hauptmann_deep_2020}~\cite{grohl_deep_2021}, con metas de
resolución espacial de 150--200~$\mu$m y profundidad de 1--2~cm en tejido articular.
El esfuerzo de instrumentación ha progresado de forma incremental: en el segundo
semestre se validó la fuente de excitación (driver en modo avalancha, láser LED de
\SI{445}{\nano\meter}, pulsos $<$\SI{30}{\nano\second})~\cite{wu_advanced_2021}
~\cite{li_photoacoustic_2020}; en el tercero se evolucionó dicha fuente (Driver V2.0)
y se validó el núcleo digital con señales sintéticas (\textit{Virtual ADC}).

El \textbf{avance central de este periodo} es llevar esa cadena al hardware físico
completo: integrar el ADC con la FPGA y validar, con señal real, un DAQ de alta
velocidad operativo. La detección ultrasónica se realizó en una primera
caracterización (Sección~\ref{sec:futuro}); su dificultad no es una valoración a
priori sino consecuencia de cuatro condiciones objetivas: la señal acústica esperada
es de microvoltios a milivoltios (cerca del piso de ruido), el transductor es una
fuente de alta impedancia, la excitación de alto voltaje acopla ruido
electromagnético, y el evento es de nanosegundos.

\section{Objetivos del Periodo}
\begin{enumerate}
    \item \textbf{Integración física ADC--FPGA:} montar el módulo AD9226 y certificar,
    con señal real, la integridad, linealidad y respuesta dinámica de la cadena de
    adquisición a la velocidad y resolución de diseño.
    \item \textbf{Operación del Driver V2.0:} verificar la estabilidad de la fuente de
    excitación y estimar la potencia de excitación.
    \item \textbf{Caracterización inicial de la detección:} realizar el primer intento
    de adquisición y discriminar rigurosamente el origen de la señal observada.
\end{enumerate}

\section{Materiales}
El sistema se construyó con componentes accesibles (Tabla~\ref{tab:materiales}). El
código de programación de la FPGA está disponible en el repositorio del
proyecto:\footnote{\url{https://github.com/zaned897/photoacoustic_daq}}

\begin{table}[htbp]
\centering\small
\begin{tabular}{ll}
\toprule
\textbf{Componente} & \textbf{Detalle} \\
\midrule
FPGA & Tang Nano 9K (Gowin GW1NR-9C), cristal 27 MHz \\
Conversor ADC & Módulo AD9226 (12 bit, hasta 65 MS/s) \\
Fuente de excitación & Driver de corriente V2.0 (validado), modo avalancha \\
Láser & Diodo LED 445 nm (Nichia NUBM44/47) \\
Transductor & Piezoeléctrico Yushi 2.5 MHz \\
Sensado de corriente & Resistencia \textit{shunt} de 0.1 $\Omega$ \\
Referencia de medición & Osciloscopio (200 MS/s) \\
\bottomrule
\end{tabular}
\caption{Materiales principales del sistema de adquisición fotoacústico.}
\label{tab:materiales}
\end{table}

\section{Arquitectura del Sistema de Adquisición (DAQ)}

El experimento completo sigue el flujo de la Figura~\ref{fig:sistema}: excitación
láser pulsada, generación fotoacústica en el tejido, adquisición en ráfaga de alta
velocidad y transmisión en el tiempo muerto entre disparos hacia la PC. Este reporte se
centra en el subsistema de adquisición (DAQ).

\figph{sistema-experimento.png}{Visión general del experimento (concepto
\textit{store-and-forward}): excitación láser $\to$ generación fotoacústica $\to$
adquisición en ráfaga $\to$ \textit{buffer} $\to$ transmisión en el tiempo muerto $\to$
procesamiento en PC. Diagrama conceptual; la implementación de este periodo opera a
54~MS/s con 270 muestras (ventana 5~µs), ver Tabla~\ref{tab:bringup}.}{sistema}

El DAQ está organizado en seis etapas encadenadas (Figura~\ref{fig:daq-arq}). El
diseño explota el bajo ciclo de trabajo del láser mediante una estrategia
\textbf{\textit{store-and-forward}}: se captura una ráfaga a alta velocidad en la
memoria interna (BRAM) y se transmite durante el tiempo muerto entre disparos,
desacoplando la velocidad de muestreo de la de transmisión.

\begin{figure}[htbp]
\centering
\resizebox{\textwidth}{!}{%
\begin{tikzpicture}[
  font=\sffamily,
  block/.style={rectangle, draw=blue!55, thick, rounded corners, fill=blue!5,
                align=center, text width=3.1cm, minimum height=2.3cm, inner sep=4pt},
  arr/.style={-{Stealth[length=3mm]}, very thick, draw=blue!55}
]
\node[block] (exc) {\textbf{1. Excitación}\\[3pt]\scriptsize Driver V2.0\\
  \scriptsize trigger 5\,kHz\\ \scriptsize monitor óptico (TIA)};
\node[block, right=0.7cm of exc] (trg) {\textbf{2. Trigger}\\[3pt]\scriptsize pin 77 / botón S2\\
  \scriptsize detección de flanco\\ \scriptsize arranque de ventana};
\node[block, right=0.7cm of trg] (adc) {\textbf{3. Conversión}\\[3pt]\scriptsize AD9226 $\cdot$ 12 bit\\
  \scriptsize 54 MS/s (PLL $\times$2)\\ \scriptsize $\pm$5 V, LSB 2.44 mV};
\node[block, right=0.7cm of adc] (cap) {\textbf{4. Captura}\\[3pt]\scriptsize FSM store-and-forward\\
  \scriptsize BRAM 270 muestras\\ \scriptsize ventana 5\,$\mu$s};
\node[block, right=0.7cm of cap] (uart) {\textbf{5. Enlace}\\[3pt]\scriptsize UART\\
  \scriptsize header sync 4 B\\ \scriptsize FW version};
\node[block, right=0.7cm of uart] (host) {\textbf{6. Host}\\[3pt]\scriptsize Python\\
  \scriptsize visualización\\ \scriptsize promediado coherente};
\draw[arr] (exc)--(trg); \draw[arr] (trg)--(adc); \draw[arr] (adc)--(cap);
\draw[arr] (cap)--(uart); \draw[arr] (uart)--(host);
\end{tikzpicture}%
}
\caption{Arquitectura del DAQ por etapas y subetapas, validada de forma incremental
(\textit{bring-up}, Tabla~\ref{tab:bringup}).}
\label{fig:daq-arq}
\end{figure}

\paragraph{Validación incremental.} Se siguió una metodología de \textit{bring-up} por
capas: cada etapa prueba una sola variable y no se avanza hasta que la anterior se
comporta exactamente como se espera (Tabla~\ref{tab:bringup}), lo que permite atribuir
cualquier fallo a lo que la etapa en curso agregó.

\begin{table}[htbp]
\centering\small
\begin{tabular}{cll}
\toprule
\textbf{Etapa} & \textbf{Subetapa que valida} & \textbf{Estado} \\
\midrule
0 & Reloj + configuración del bitstream & \checkmark \\
1 & Trigger (botón S2 / pin 77) y detección de flanco & \checkmark \\
2 & Enlace UART aislado (contador) & \checkmark \\
3 & Lectura del ADC en vivo (voltímetro) & \checkmark \\
4 & Captura por reloj interno (5 kHz) & \checkmark \\
5 & Captura por trigger externo & \checkmark \\
6 & Ráfaga de 270 muestras @ 27 MS/s & \checkmark \\
7 & Ráfaga @ 54 MS/s (PLL $\times$2) & \checkmark \\
8 & Ráfaga @ 54 MS/s + 12 bits (LSB 2.44 mV) & \checkmark \\
\bottomrule
\end{tabular}
\caption{Validación incremental (\textit{bring-up}) del DAQ.}
\label{tab:bringup}
\end{table}

\subsection{Algoritmo de captura y sincronización}
El núcleo del DAQ es una máquina de estados (FSM) \textit{store-and-forward}
(Figura~\ref{fig:fsm}): en reposo (IDLE) vigila el trigger; al detectar un flanco pasa
a CAPTURE y escribe $N$ muestras en la BRAM al ritmo del reloj de muestreo; luego
emite un encabezado de sincronía (PREP, HDR) y transmite por UART (SEND) antes de
volver a IDLE. Fuera de IDLE ignora nuevos triggers, garantizando ráfagas íntegras.

\begin{figure}[htbp]
\centering
\resizebox{0.9\textwidth}{!}{%
\begin{tikzpicture}[
  font=\sffamily,
  st/.style={rectangle, draw=blue!55, thick, rounded corners, fill=blue!5,
             minimum width=1.8cm, minimum height=0.9cm, align=center},
  tr/.style={-{Stealth[length=2.5mm]}, thick, draw=blue!55}
]
\node[st] (idle) {IDLE};
\node[st, right=1.4cm of idle] (cap) {CAPTURE};
\node[st, right=1.4cm of cap] (prep) {PREP};
\node[st, right=1.4cm of prep] (hdr) {HDR};
\node[st, right=1.4cm of hdr] (send) {SEND};
\draw[tr] (idle) -- node[above,font=\scriptsize]{trig} (cap);
\draw[tr] (cap) -- node[above,font=\scriptsize]{ptr=N-1} (prep);
\draw[tr] (prep) -- (hdr);
\draw[tr] (hdr) -- node[above,font=\scriptsize]{hdr$\times$4} (send);
\draw[tr] (send) to[out=-90,in=-90] node[below,font=\scriptsize]{ptr=N} (idle);
\end{tikzpicture}%
}
\caption{Máquina de estados del núcleo de adquisición (\textit{store-and-forward}).}
\label{fig:fsm}
\end{figure}

\subsection{El costo del ciclo de lectura completo}
El cristal de 27~MHz se duplica a \textbf{54~MS/s} mediante el PLL interno (rPLL,
$\times2$; periodo de muestreo de 37 a 18.5~ns), y se pasó de 8 a \textbf{12 bits}
reales del AD9226 (de 256 a 4096 niveles, LSB de 39 a \textbf{2.44~mV};
Figura~\ref{fig:niveles}). Para el transductor de 2.5~MHz esto da $\sim$21.6
muestras/ciclo, holgado sobre Nyquist, y cada punto de adquisición es una muestra de
12 bits cada 18.5~ns (Figura~\ref{fig:daqmuestreo}).

Llevar el ciclo de lectura a esta resolución y tasa multiplica el volumen de datos por
ventana ($N\times b$ bits), con consecuencias en tres frentes:
\begin{itemize}
    \item \textbf{Memoria (BRAM):} cada ráfaga ocupa $N\times b$ bits; aun una ventana
    de profundidad útil ($\sim$5.8~cm, 37.7~µs $\to$ $\sim$24~Kbit a 12 bits) cabe
    holgadamente en el BSRAM del GW1NR-9C ($\approx$468~Kbit), por lo que la memoria
    \emph{no} es el factor limitante (Figura~\ref{fig:memoria}).
    \item \textbf{Transmisión UART:} más bits por muestra y más muestras alargan la
    transmisión por ráfaga; como esta excede el periodo del trigger (200~µs a 5~kHz),
    la tasa efectiva de ráfagas baja y se recurre al \textbf{promediado coherente}
    ($\mathrm{SNR}\propto\sqrt{N}$) en lugar de capturar cada disparo
    (Figura~\ref{fig:timing}).
    \item \textbf{Sincronización de la ventana:} la captura arranca en el flanco del
    trigger, sin pre-trigger; al duplicar la tasa, la ventana temporal de un número
    fijo de muestras se acorta (270 muestras: 10~µs a 27~MS/s $\to$ 5~µs a 54~MS/s),
    lo que exige mayor precisión en la sincronía disparo–ventana.
\end{itemize}

\figph{fig-timing-fsm.png}{Sincronización \textit{store-and-forward}: la captura
(5~µs) es muy breve frente a la transmisión por UART (varios ms), durante la cual se
omiten disparos.}{timing}

\figph{fig-niveles-codificacion.png}{Profundidad de codificación: a 12 bits el paso de
cuantización (LSB 2.44~mV) resuelve detalles que a 8 bits (LSB 39~mV) se
pierden.}{niveles}

\figph{daq-muestreo.png}{Muestreo real del DAQ. Arriba: ventana de 270 muestras (cada
punto = una muestra). Abajo: zoom al transitorio (\textit{stem}); el evento queda en
$\sim$1 muestra.}{daqmuestreo}

\figph{fig-memoria.png}{Memoria por ráfaga frente a la ventana de captura (8/10/12
bits a 54~MS/s). El BSRAM del FPGA excede ampliamente lo necesario.}{memoria}

\section{Metodología y Resultados}
\label{sec:resultados}

\subsection{Integración y validación del DAQ con ADC físico}
El módulo AD9226 se integró con la FPGA sobre una placa de acrílico que aloja,
atornilladas, las PCBs de todas las etapas (Figura~\ref{fig:acrilico}); la etapa de
adquisición se verificó por partes —trigger, conversión, captura y transmisión—
antes de integrarla (Figura~\ref{fig:fpga-etapas}). Con el ADC físico operando, la
cadena \textbf{captura, almacena y transmite sin pérdida de datos} a 54~MS/s con 12
bits (LSB $\approx \SI{2.44}{\milli\volt}$); la validación por etapas
(Figura~\ref{fig:bringup}) y la captura de una ráfaga (Figura~\ref{fig:captura})
confirman la velocidad y los niveles de diseño. La adquisición y visualización en el
host se realizan con software propio (Figura~\ref{fig:software}).

\figph{placa-acrilico.png}{Plataforma de integración: placa de acrílico con las PCBs
de todas las etapas atornilladas.}{acrilico}

\figph{daq-fpga-etapas.png}{Etapa de adquisición del FPGA: hardware (FPGA + AD9226) y
diagrama de bloques (trigger, conversión, captura en BRAM, UART).}{fpga-etapas}

\figph{bringup-etapas.png}{Validación incremental del FPGA por etapas (0--8).}{bringup}

\figph{daq-captura.png}{Captura del DAQ (corriente del diodo vía \textit{shunt} de
0.1~$\Omega$): ventana de 270 muestras @ 54~MS/s, 12 bit; una adquisición y su
promedio coherente $\times$32.}{captura}

\figph{daq-python-software.png}{Software final de adquisición y
visualización (GUI): ventana completa, zoom de región y histórico, con promedio
coherente y lectura en mV/código.}{software}

\paragraph{Comparación con la referencia.} Para validar la adquisición sobre una señal
real rápida se capturó el mismo transitorio sincronizado al láser con el osciloscopio
de referencia (200~MS/s) y con el DAQ (54~MS/s, 12~bit). La Figura~\ref{fig:daqvsosc},
con ambas trazas alineadas en el transitorio, muestra que el DAQ reproduce el evento
con \textit{timing} coherente. Ambos miden el mismo disparo en magnitudes distintas:
el osciloscopio la tensión en el láser y el DAQ la corriente del diodo en el
\textit{shunt}. A 54~MS/s el transitorio ($\sim$25~ns FWHM en la referencia) queda en
$\sim$1 muestra: el DAQ capta el evento y su amplitud, no su estructura de
nanosegundos. Dicho transitorio corresponde a la descarga de avalancha que excita el
diodo (técnica de conmutación en avalancha de transistores~\cite{williams_an47}), de
naturaleza ultracorta.

\figph{daq-vs-osc.png}{Mismo disparo medido por dos vías: tensión en el láser
(osciloscopio, 200~MS/s) y corriente del diodo en el \textit{shunt} (DAQ, 54~MS/s),
alineadas en el evento. \textit{Timing} coherente; magnitudes verticales distintas
(tensión vs corriente).}{daqvsosc}

\subsection{Driver de excitación V2.0 y estimación de potencia}
El Driver V2.0 entregó pulsos de corriente estables, registrados como la caída de
tensión en el \textit{shunt} de 0.1~$\Omega$ y digitalizados por el DAQ
(Figura~\ref{fig:pulsos}). Combinando esa corriente con la tensión del diodo medida en
el osciloscopio se obtuvo una \textbf{estimación preliminar} de la potencia eléctrica
de excitación (Figura~\ref{fig:potencia}), de orden de magnitud: ambas magnitudes
provienen de adquisiciones no simultáneas del disparo —periódico y reproducible— y la
corriente está submuestreada (cota inferior del pico).

\figph{pulsos-corriente.png}{Pulso de corriente del diodo (caída en el \textit{shunt}
de 0.1~$\Omega$ digitalizada por el DAQ).}{pulsos}

\figph{potencia-laser.png}{Estimación preliminar de la potencia de excitación a partir
de la tensión del diodo (osciloscopio) y la corriente del \textit{shunt} (DAQ),
alineadas en el transitorio. Valor de orden de magnitud.}{potencia}

\subsection{Caracterización inicial de la detección}
En el osciloscopio se observó una señal en el transductor con la apariencia de un
\textit{ringing} acústico ($\sim$2.5~MHz, $\sim$150~mV pp), coincidente con el disparo
(Figura~\ref{fig:pickup}). Dos hechos la descartan como detección acústica:

\begin{enumerate}
    \item \textbf{Es acoplamiento electromagnético (\textit{pickup}).} Tres controles
    lo demuestran: la onda persiste con el transductor \emph{desconectado}; coincide
    con el disparo en $t=0$ (sin retardo de tiempo de vuelo); y no cambia al variar la
    distancia fuente--sensor.
    \item \textbf{Se pierde al pasar al ADC.} La señal solo era visible en el
    osciloscopio (entrada de 1~M$\Omega$): al conectar el transductor al ADC, se
    perdía. Esto evidencia el acoplamiento de impedancias (Sección~\ref{sec:disc-imp}):
    la entrada del digitalizador carga la fuente de alta impedancia del piezo.
\end{enumerate}

Una prueba con fuente mecánica conocida (rotura de mina, Hsu--Nielsen, sin láser)
tampoco produjo señal observable, consistente con la ausencia de pre-amplificación y
con que dicha fuente concentra su energía por debajo de \SI{1}{\mega\hertz}, mal
acoplada a un transductor de 2.5~MHz. La caracterización localiza con precisión la
etapa pendiente y orienta el trabajo futuro (Sección~\ref{sec:futuro}).

\figph{ringing-falso.png}{Control de descarte (osciloscopio). CH1: pulso láser; CH2:
onda de $\sim$150~mV pp ($\sim$2.6~MHz) coincidente con el disparo. Persiste con el
sensor desconectado y llega en $t=0$: es \textit{pickup}, no detección
acústica.}{pickup}

\section{Discusión}

\subsection{El DAQ como contribución central}
El resultado principal es un \textbf{DAQ de alta velocidad, validado y de bajo costo}:
captura sin pérdida a 54~MS/s con 12 bits, un sobremuestreo holgado respecto a la
frecuencia central del transductor (2.5~MHz). La arquitectura \textit{store-and-forward}
sostiene ese muestreo sobre un enlace serial sencillo y la validación por capas
certifica cada etapa de forma independiente. Esta plataforma es la base instrumental
sobre la cual se atacará la detección.

\subsection{La detección y el acoplamiento de impedancias}
\label{sec:disc-imp}
Los controles separan con nitidez los subsistemas ya validados (ADC, FPGA, enlace,
Driver V2.0) del eslabón pendiente, la interfaz analógica de detección. El obstáculo
clave, además del piso de ruido y la EMI, es el \textbf{acoplamiento de impedancias}:
el piezo es una fuente de alta impedancia y conectarlo a una entrada de baja impedancia
forma un divisor que atenúa la señal —precisamente lo observado, pues la señal visible
en el osciloscopio (1~M$\Omega$) se perdía al pasar al ADC. La
Figura~\ref{fig:impedancia} cuantifica este obstáculo: una entrada de 50~$\Omega$
retiene apenas una fracción, mientras que una de $\sim$1~M$\Omega$ la preserva. Su
solución —y por tanto el primer requisito de diseño del front-end— es un
\textbf{pre-amplificador de alta impedancia de entrada junto al sensor}: además de
elevar la señal ($\mu$V--mV) por encima del piso de observación, su salida de baja
impedancia transporta la señal sin atenuación ni susceptibilidad al \textit{pickup}.

\figph{fig-impedancia.png}{El acoplamiento de impedancias como obstáculo de diseño:
fracción de señal retenida frente a la impedancia de entrada del receptor, para dos
impedancias de fuente del piezo. Una entrada de baja impedancia (50~$\Omega$, o la del
ADC) atenúa severamente —de ahí que la señal se perdiera al pasar del osciloscopio
(1~M$\Omega$) al ADC—; el pre-amplificador de alta-Z resuelve este
acoplamiento.}{impedancia}

\section{Limitaciones}
\begin{itemize}
    \item La detección acústica aún no se demuestra por encima del piso de observación;
    es el objeto del trabajo futuro.
    \item La ventana de captura actual (5~$\mu$s, $\approx \SI{7.7}{\milli\meter}$) es
    insuficiente para la cobertura de tejido objetivo; su ampliación requiere
    \textit{buffers} (FIFO) en la FPGA.
    \item La estimación de potencia es de orden de magnitud (capturas no simultáneas,
    corriente submuestreada).
\end{itemize}

\section{Conclusiones}
\begin{enumerate}
    \item \textbf{DAQ validado (avance central):} cadena con ADC físico certificada
    —12 bits, 54~MS/s, sin pérdida de datos— a la velocidad y niveles de diseño.
    \item \textbf{Driver V2.0 y potencia:} operación estable y estimación preliminar de
    la potencia de excitación.
    \item \textbf{Detección localizada:} la señal observada es \textit{pickup} (no
    acústica) y se pierde al pasar al ADC, acotando la etapa crítica al front-end
    analógico (acoplamiento de impedancias + ruido).
\end{enumerate}

Aunque el objetivo del periodo —la validación completa de la cadena fotoacústica, con
el transductor excitado y sensando— no se alcanzó en su totalidad, haber diseñado e
implementado una FPGA como sistema de adquisición capaz de operar a las frecuencias de
muestreo y resolución requeridas, y con materiales de bajo costo, constituye un
\textbf{hito} del proyecto: la plataforma de adquisición queda resuelta y la detección
deja de ser una incógnita difusa para convertirse en un problema acotado de front-end
analógico.

\section{Trabajo Futuro}
\label{sec:futuro}

El foco del siguiente periodo es la cadena de detección, sobre la base del DAQ ya
validado. Se validará por capas, en analogía con el \textit{bring-up} digital:
salud eléctrica del piezo (capacitancia/continuidad) $\to$ respuesta mecánica directa
sin láser $\to$ fuente acústica conocida acoplada (\textit{in-band}) $\to$
incorporación del pre-amplificador $\to$ reintroducción del láser buscando la señal en
la ventana retardada por tiempo de vuelo.

El \textbf{pre-amplificador} es el elemento decisivo: alta impedancia de entrada
($\geq \SI{1}{\mega\ohm}$), ganancia 40--60~dB, ancho de banda 0.1--5~MHz, ruido
$\leq$ 3--5~nV/$\sqrt{\text{Hz}}$ y apantallamiento (Tabla~\ref{tab:preamp}). En
paralelo se controlará la EMI (sincronía por disparo óptico, apantallamiento, ventana
retardada/pre-trigger) y se ampliará la ventana de captura con \textit{buffers} (FIFO)
hacia profundidades útiles, para luego adquirir A-scans en fantomas (gelatina +
grafito), reconstruir por retroproyección y aplicar una primera etapa de IA de
reducción de ruido. Sin \textit{pulser} ni segundo transductor, la validación acústica
se prioriza por viabilidad (láser sobre absorbente fuerte; pulso-eco con generador).

\begin{table}[htbp]
\centering\small
\begin{tabular}{lll}
\toprule
\textbf{Candidato} & \textbf{Tipo} & \textbf{Ruido} \\
\midrule
AD8331/AD8332 & LNA + VGA (ultrasonido) & $\sim$0.74 nV/$\sqrt{\text{Hz}}$; opción primaria \\
OPA657 & Op-amp JFET & 4.8 nV/$\sqrt{\text{Hz}}$, GBW 1.6 GHz \\
AD8067 + JFET (BF862) & Amplificador de carga & robusto a la capacitancia del cable \\
\bottomrule
\end{tabular}
\caption{Candidatos a preamplificador (opción primaria: AD8331/AD8332).}
\label{tab:preamp}
\end{table}

\begin{table}[htbp]
\centering\small
\begin{tabular}{|p{0.13\textwidth}|p{0.42\textwidth}|p{0.34\textwidth}|}
\hline
\textbf{Bloque} & \textbf{Actividad} & \textbf{Entregable} \\ \hline
B1 & Diseño del front-end (preamp, impedancias, apantallamiento, sincronía) &
Esquemático + presupuesto de ruido \\ \hline
B2 & Fabricación y validación de la recepción & Cadena caracterizada \\ \hline
B3 & Firmware: ventana ampliada (FIFO) + pre-trigger & A-scan de profundidad útil \\ \hline
B4 & Adquisición PA en fantoma & A-scan con tiempo de vuelo coherente \\ \hline
B5 & Procesamiento (retroproyección) + IA (denoising) & Reconstrucción y artículo \\ \hline
\end{tabular}
\caption{Cronograma propuesto para el siguiente periodo (jul.--dic.\ 2026).}
\label{tab:cronograma_rev}
\end{table}

La metodología de validación por capas permitió afirmar con certeza que el subsistema
digital está resuelto y señalar el analógico como el reto pendiente: ese aislamiento es,
en sí mismo, parte del avance. La lección —que en sistemas fotoacústicos de bajo costo
el front-end analógico y la integridad electromagnética dominan el éxito tanto como el
cómputo— es una contribución relevante para el artículo de instrumentación previsto.

\printbibliography

\end{document}
